% Part: incompleteness
% Chapter: arithmetization-syntax
% Section: coding-symbols

\documentclass[../../../include/open-logic-section]{subfiles}

\begin{document}

\olfileid{inc}{art}{cod}
\olsection{ચિહ્નોનું સંકેતબદ્ધીકરણ}

પ્રથમ-ક્રમના તર્કશાસ્ત્રની મૂળભૂત ભાષા~$\Lang L$
નીચેના ચિહ્નો વાપરે છે:
\[
\lfalse \quad \lnot \quad \lor \quad \land \quad \lif \quad \lforall
\quad \lexists \quad \eq \quad ( \quad ) \quad ,
\]
આ સાથે ચલો અને !!{constant}sના
!!{enumerable} ગણો તથા મનસ્વી સ્થાનિકતાવાળા
!!{function}s અને !!{predicate}sના
!!{enumerable} ગણો હોય છે. આપણે દરેક
ચિહ્નને \emph{સંકેત નંબર} એવી રીતે આપી
શકીએ કે દરેક ચિહ્નને અનન્ય નંબર મળે અને
બે જુદાં ચિહ્નોને ક્યારેય એ જ નંબર ન મળે.
આ શક્ય છે કારણ કે બધા ચિહ્નોનો ગણ
!!{enumerable} છે અને તેથી તેની તથા
પ્રાકૃતિક સંખ્યાઓના ગણની વચ્ચે
!!a{bijection} છે. પરંતુ સંકેત નંબરમાંથી
ચિહ્ન તથા તેના વિશેની કેટલીક માહિતી
(દાખલા તરીકે, !!a{function}ની સ્થાનિકતા)
સંગણનીય રીતે પાછી મેળવી શકાય તેની
ખાતરી પણ જોઈએ. એ કરવાની ઘણી રીત છે.
અહીં આદિમ પુનરાવર્તી વિધેયો વાપરતી
એક રીત આપી છે. (યાદ કરો કે
$\tuple{n_0, \dots, n_k}$ એ સંખ્યાઓ
$n_0$, \dots, $n_k$ના અનુક્રમનો સંકેત
નંબર છે.)

\begin{defn}
$s$ એ~$\Lang L$નું ચિહ્ન હોય તો
તેનો \emph{ચિહ્ન-સંકેત}~$\scode s$
આ પ્રમાણે વ્યાખ્યાયિત કરીએ:
\begin{enumerate}
\item $s$ તાર્કિક ચિહ્નોમાંનું હોય તો
  $\scode s$ નીચેના કોષ્ટક મુજબ છે:
\[
\begin{array}{cccccccccc}
  \lfalse & \lnot & \lor & \land & \lif & \lforall \\
\tuple{0, 0} & \tuple{0, 1} & \tuple{0, 2} & \tuple{0, 3} &
\tuple{0, 4} & \tuple{0, 5} \\
\lexists & \eq & ( & ) & ,\\
\tuple{0, 6} & \tuple{0, 7} &
\tuple{0, 8} & \tuple{0, 9} & \tuple{0, 10}
\end{array}
\]
\item $s$ એ~$i$મો ચલ~$\Obj v_i$ હોય, તો $\scode s = \tuple{1, i}$.
\item $s$ એ~$i$મો !!{constant}~$\Obj c_i$ હોય, તો
  $\scode s = \tuple{2, i}$.
\item $s$ એ~$i$મો $n$-સ્થાનિક !!{function}~$\Obj f_i^n$ હોય, તો
  $\scode s = \tuple{3, n, i}$.
\item $s$ એ~$i$મો $n$-સ્થાનિક !!{predicate}~$\Obj P_i^n$ હોય, તો
  $\scode s = \tuple{4, n, i}$.
\end{enumerate}
\end{defn}

\begin{prop}
નીચેના સંબંધો આદિમ પુનરાવર્તી છે:
\begin{enumerate}
\item $\fn{Fn}(x, n)$ ત્યારે અને માત્ર ત્યારે,
  જ્યારે કોઈ~$i$ માટે~$x$ એ~$\Obj f^n_i$નો
  સંકેત નંબર હોય; એટલે~$x$ એ $n$-સ્થાનિક
  !!{function}નો સંકેત નંબર હોય.
\item $\fn{Pred}(x, n)$ ત્યારે અને માત્ર ત્યારે,
  જ્યારે કોઈ~$i$ માટે~$x$ એ~$\Obj P^n_i$નો
  સંકેત નંબર હોય અથવા~$x$ એ~$\eq$નો
  સંકેત નંબર અને~$n = 2$ હોય; એટલે~$x$
  એ $n$-સ્થાનિક !!{predicate}નો સંકેત
  નંબર હોય.
\end{enumerate}
\end{prop}

\begin{defn}
$s_0, \dots, s_{n-1}$ ચિહ્નોનો અનુક્રમ
હોય, તો તેની \emph{ગડેલ-સંખ્યા}
$\tuple{\scode{s_0}, \dots, \scode{s_{n-1}}}$ છે.
\end{defn}

\begin{explain}
\emph{સંકેત નંબરો} અને \emph{ગડેલ-સંખ્યાઓ}
અલગ વસ્તુઓ છે. દાખલા તરીકે, ચલ~$\Obj v_5$
નો સંકેત નંબર~$\scode{\Obj
  v_5} = \tuple{1, 5} = 2^2\cdot 3^6$ છે.
પરંતુ પદ તરીકે વિચારેલું ચલ~$\Obj v_5$
લંબાઈ~$1$ ધરાવતો ચિહ્નોનો અનુક્રમ પણ
છે. \emph{પદ}~$\Obj v_5$ની
\emph{ગડેલ-સંખ્યા}~$\Gn{\Obj v_5}$ એ
$\tuple{\scode{\Obj v_5}} = 2^{\scode{\Obj
    v_5} + 1} = 2^{2^2\cdot 3^6 + 1}$ છે.
\end{explain}

\begin{ex}
યાદ કરો કે $k_0$, \dots, $k_{n-1}$
સંખ્યાઓનો અનુક્રમ હોય, તો અવિભાજ્ય
સંખ્યાઓની ઘાતો વડે થતા સંકેતબદ્ધીકરણમાં
અનુક્રમ~$\tuple{k_0, \dots, k_{n-1}}$નો
સંકેત નંબર આ છે:
\[
2^{k_0+1}\cdot3^{k_1+1}\cdot \dots \cdot p_{n-1}^{k_{n-1}+1},
\]
અહીં~$p_i$ એ~$i$મી અવિભાજ્ય સંખ્યા છે
(શરૂઆત~$p_0 = 2$થી થાય છે). તેથી,
દાખલા તરીકે, સૂત્ર~$\eq[\Obj v_0][\Obj 0]$
અથવા વધુ સ્પષ્ટ રીતે~${\eq}(\Obj v_0,\Obj c_0)$ની
ગડેલ-સંખ્યા છે:
\[
\tuple{\scode{\eq},\scode{(},\scode{\Obj v_0},\scode{,}, \scode{\Obj
    c_0},\scode{)}}.
\]
અહીં $\scode{\eq}$ એ~$\tuple{0,7} = 2^{0+1}\cdot
3^{7+1}$ અને $\scode{\Obj v_0}$ એ~$\tuple{1,0} = 2^{1+1}\cdot3^{0+1}$,
વગેરે. તેથી~$\Gn{=(\Obj v_0,\Obj c_0)}$ એ
\begin{multline*}
2^{\scode{=} + 1}\cdot 3^{\scode{(}+1}\cdot 5^{\scode{\Obj v_0}+1}
\cdot 7^{\scode{,} + 1} \cdot 11^{\scode{\Obj c_0}+1} \cdot
13^{\scode{)}+1} = \\
2^{2^1\cdot 3^8 + 1}\cdot 3^{2^1\cdot 3^9+1}\cdot 5^{2^2\cdot 3^1+1}
\cdot 7^{2^1\cdot 3^{11} + 1} \cdot 11^{2^3\cdot3^1+1} \cdot
13^{2^1\cdot3^{10}+1} = \\
2^{13\,123}\cdot 3^{39\,367}\cdot 5^{13}\cdot 7^{354\,295}\cdot11^{25}\cdot13^{118\,099}.
\end{multline*}
\end{ex}

\end{document}
